\documentclass[10pt,a4paper]{amsart}
\usepackage[margin=19mm]{geometry}
\usepackage{amsmath,amssymb,mathtools}
\usepackage[scr=rsfs,cal=euler]{mathalfa}
\usepackage{xfrac}
\usepackage[unicode,hidelinks,bookmarks=false]{hyperref}
\newcommand{\ie}{i.e.\ }
\newcommand{\cf}{cf.\ }
\newcommand{\resp}{respectively\ }
\newcommand{\Subsection}{\subsection}
\providecommand{\nobreakdash}{\nobreak\hbox{-}}
\setlength{\emergencystretch}{2em}
\multlinegap0pt
\usepackage{float}
\usepackage{amssymb}
\usepackage{xcolor}
\newtheorem{thm}{Theorem}
\title{Stability Regions and Bifurcations for Higher-Order Fractional Difference Equations}
\author{Janardhan Chevala, Sachin Bhalekar}
\date{Source: arXiv:2603.23090v2 (CC BY 4.0)}
\begin{document}
\maketitle

\noindent\emph{Source and licence.} Stability Regions and Bifurcations for Higher-Order Fractional Difference Equations. Version arXiv:2603.23090v2 (CC BY 4.0), distributed by arXiv under the Creative Commons Attribution 4.0 International licence, http://creativecommons.org/licenses/by/4.0/, as recorded in the arXiv licence metadata for this exact version and shown on the abstract page https://arxiv.org/abs/2603.23090v2. Full licence notice: \texttt{licenses/arxiv-2603.23090-CC-BY-4.0.txt}.\par\smallskip

\noindent\textit{Editorial scope.} Counted unit: the article's thm thm:bifurcation-values (source0.tex lines 354--368). The article's complete original proof of it is delivered with it (source0.tex lines 370--400, one \texttt{proof} environment of 31 lines). No mathematics is written, repaired or re-derived: the statement and the proof are the article's own lines, in the article's own notation, and no retained line is substituted or re-flowed.  Retained uncounted as the source-local context the proof needs: lines 211--216.  Nothing was omitted from any retained span, so the package carries no omission record.  No retained line contains a citation, so the wrapper carries no bibliography and no bibliography entry.  No retained line contains a figure environment, an image-inclusion command or drawing code, so no image is delivered and the package declares no asset dependency.  Editorial formatting only, in addition: an AMS \texttt{amsart} preamble that restates the article's own declarations and macros used by the retained lines, namely the environments \texttt{thm} on separate counters, together with the standard packages those lines need.  The retained environments print in this document as Theorem 1 in order of appearance, whereas the article prints the same items with the numbering of its own full document.  The complete native source of the article as distributed in the arXiv e-print for this exact version is bundled unchanged as \texttt{sources/197052/source0.tex}.
 \begin{align}
 b ={}& 2^\alpha\left(\sin\frac{\theta}{2}\right)^\alpha
 e^{\mathrm{i}\left[\frac{\alpha\pi}{2}+\theta\left(2-\frac{\alpha}{2}\right)\right]} \nonumber\\
 &+ a\,2^\beta\left(\sin\frac{\theta}{2}\right)^\beta
 e^{\mathrm{i}\left[\frac{\beta\pi}{2}+\theta\left(1-\frac{\beta}{2}\right)\right]} + 1. \label{simbval}
\end{align}

\begin{thm}[Boundary bifurcation values]\label{thm:bifurcation-values}
Let $0 < \beta \leq 1 < \alpha \leq 2$ and let $\gamma$ be given by \eqref{simbval}. Define
\begin{equation}\label{a1a2}
a_1=2^{\alpha-\beta},
\qquad
a_2=2^{\alpha-\beta}\frac{4-\alpha}{2-\beta}.
\end{equation}
Then:
\begin{enumerate}
\item $\gamma(0)=\gamma(\pi)$ if and only if $a=a_1$;
\item $\gamma'(\pi)=0$ if and only if $a=a_2$;
\item $0<a_1<a_2$.
\end{enumerate}
Thus, $a_1$ is the unique value for which $\gamma(0)=\gamma(\pi)$, and $a_2$ is the unique value for which the boundary is stationary at $\theta=\pi$.
\end{thm}

\begin{proof}
Since $\sin(0)=0$, equation \eqref{simbval} gives $\gamma(0)=1$. At $\theta=\pi$, the two phase factors satisfy
\begin{equation*}
e^{\mathrm{i}[\alpha\pi/2+\pi(2-\alpha/2)]}=1,
\qquad
e^{\mathrm{i}[\beta\pi/2+\pi(1-\beta/2)]}=-1.
\end{equation*}
Therefore,
\begin{equation*}
\gamma(\pi)=1+2^\alpha-a\,2^\beta.
\end{equation*}
The equality $\gamma(0)=\gamma(\pi)$ is equivalent to $2^\alpha=a\,2^\beta$, and hence to $a=a_1$.

To obtain the second value, differentiate \eqref{simbval}. The derivatives of the radial factors vanish at $\theta=\pi$ because $\cos(\pi/2)=0$. Consequently,
\begin{align*}
\gamma'(\pi)
&=\mathrm{i}\,2^\alpha\left(2-\frac{\alpha}{2}\right)
-\mathrm{i}\,a\,2^\beta\left(1-\frac{\beta}{2}\right)\\
&=\mathrm{i}\left[2^{\alpha-1}(4-\alpha)
-a\,2^{\beta-1}(2-\beta)\right].
\end{align*}
Thus, $\gamma'(\pi)=0$ if and only if
\begin{equation*}
a=2^{\alpha-\beta}\frac{4-\alpha}{2-\beta}=a_2.
\end{equation*}
Finally,
\begin{equation*}
\frac{a_2}{a_1}=\frac{4-\alpha}{2-\beta}>1,
\end{equation*}
because $2-\alpha+\beta>0$. Hence $0<a_1<a_2$.
\end{proof}

\end{document}
